\section{The candidate-set sweep: width has to be trained}
\label{sec:candidates}

Every suite cell fixes its task's label width, so the suite cannot say what
happens to accuracy as the candidate set widens. That is the question the
commercial claim turns on, and \S\ref{sec:future} named a candidate-set sweep as
the direct test. The obvious implementation of that sweep is not an experiment.
With one fixed model and one fixed bank, restricting a row to a subset that still
contains the gold can only delete a wrong argmax: if the full-bank argmax is the
gold it remains the argmax of the subset, and if it is not, deleting it can only
help. Restricted accuracy is therefore at least the full-bank accuracy by
construction, whatever the model learned, and a curve produced that way would
``hold'' for a model that had learned nothing. Width has to be trained.

The sweep therefore trains a fresh decision head for each width. Every arm starts
from the same pretrained checkpoint (\SuiteCkptLarge, step \SuiteCkptLargeStep),
uses the suite's frozen fine-tuning recipe --- 4{,}000 stratified training rows,
3{,}000 optimizer steps, batch 32, LR \(10^{-4}\), AdamW, seed 0, the
checkpoint's 1{,}024-byte context --- and is tested on the suite's 2{,}000 test
rows at the width it was trained at. The only difference between arms is
\(|C|\): each training row is packed against its gold plus \(|C|-1\) distractors
drawn uniformly without replacement, with the sampling seed fixed by a hash of
task, width, step and row. Three wide-bank tasks are used: Banking77
(\(K=\EBankK\)), CLINC150 (\(K=\EClincK\)) and MASSIVE intent
(\(K=\EMassiveK\)), at \(|C| \in \{2, 8, 32, K\}\).

The design, the sampling rule and the decision rule were written and hashed
before any valid accuracy existed; both amendments, including the one that voids
the first launch, are inside the hashed file (\texttt{results/e4\_candidate\_sweep/}, SHA-256
\EPreregSha). A first launch was voided by the harness check: the width-\(K\)
Banking77 arm scored outside the 0.2-point tolerance because the decision head
was not seeded, so every arm had started from a different random
initialisation. Those runs are preserved under \texttt{\_invalid\_init} with
their reason and are excluded from every number here; the trainer now seeds
before building the head, and the design is unchanged. The pre-registered primary rule is a non-inferiority test: the
claim that accuracy holds at \(|C|=32\) survives for a task only if the paired,
one-sided 95\% bootstrap lower bound on
\(\mathrm{acc}(32)-\mathrm{acc}(K)\) is above \(-5.0\) points. The pre-registration also set a harness check: the
width-\(K\) arm, trained by the same code and schedule as the suite cell, should
reproduce the published value within 0.2 points, and any larger deviation is
reported as a discrepancy rather than repaired by substituting the published
number. Table~\ref{tab:e4} reports every arm; the full-bank arms are the same
recipe as the suite cells, not the published numbers.

\begin{table}[htbp]
\centering
\footnotesize
\caption{The candidate-set sweep at 113.47M. Each row is a separate fine-tune of
the same pretrained checkpoint, at the suite's frozen recipe (4{,}000
stratified rows, 3{,}000 steps, batch 32, LR $10^{-4}$, seed 0), differing only
in the number of candidates $|C|$ packed into every training and test row. Test
rows are the same 2{,}000 rows at every width. A fixed model restricted to a
sub-bank is not an experiment: with the gold always present, deleting candidates
can only remove wrong argmaxes, so restricted accuracy is guaranteed to be at
least the full-bank accuracy. Width is therefore trained, not simulated. Wilson
intervals are 95\%; NLL is over the sampled candidate set and is not comparable
across widths. Read from \texttt{results/e4\_candidate\_sweep/SUMMARY.json};
\texttt{scripts/e4\_summary.py} verifies every per-example record before
writing it.}
\label{tab:e4}
\setlength{\tabcolsep}{2.5pt}
\begin{tabular}{@{}lrrrrrr@{}}
\toprule
Task & $|C|$ & Accuracy & 95\% CI & NLL & $\Delta$ vs $K$ & $p$ (32 vs $K$) \\ 
\midrule
Banking77 ($K=77$) & 2 & 98.05 & 97.35--98.57 & 0.06 & $+8.90$ & --- \\
Banking77 ($K=77$) & 8 & 95.95 & 94.99--96.73 & 0.15 & $+6.80$ & --- \\
Banking77 ($K=77$) & 32 & 92.10 & 90.84--93.20 & 0.34 & $+2.95$ & $2.08\times10^{-6}$ \\
Banking77 ($K=77$) & 77 & 89.15 & 87.71--90.44 & 0.56 & --- & --- \\
CLINC150 ($K=151$) & 2 & 98.00 & 97.29--98.53 & 0.06 & $+10.30$ & --- \\
CLINC150 ($K=151$) & 8 & 95.25 & 94.23--96.10 & 0.21 & $+7.55$ & --- \\
CLINC150 ($K=151$) & 32 & 92.20 & 90.94--93.30 & 0.34 & $+4.50$ & $2.63\times10^{-12}$ \\
CLINC150 ($K=151$) & 151 & 87.70 & 86.19--89.07 & 0.60 & --- & --- \\
MASSIVE intent ($K=60$) & 2 & 96.70 & 95.82--97.40 & 0.10 & $+14.80$ & --- \\
MASSIVE intent ($K=60$) & 8 & 91.40 & 90.09--92.55 & 0.39 & $+9.50$ & --- \\
MASSIVE intent ($K=60$) & 32 & 82.85 & 81.14--84.44 & 0.91 & $+0.95$ & $2.06\times10^{-1}$ \\
MASSIVE intent ($K=60$) & 60 & 81.90 & 80.15--83.53 & 1.16 & --- & --- \\
\bottomrule
\end{tabular}
\par\smallskip
\footnotesize
$\Delta$ is in percentage points against the full-bank arm of the same task; the
$p$ column is the exact two-sided McNemar test for 32 against $K$ on the same
test rows. The full-bank arms are the harness check: the pre-registered tolerance against
the published suite cell is 0.2 points, and a larger deviation is reported as a
discrepancy rather than replaced by the published value.
\end{table}


\paragraph{What the sweep found.} On Banking77 the width-32 arm reaches
\EBankWThirtyTwo\% (95\% Wilson \EBankWThirtyTwoLo--\EBankWThirtyTwoHi) against \EBankWK\% at
width \(K\), a difference of \EBankDelta{} points with a lower bound of
\EBankLower{}; the rule's verdict is that the claim \EBankVerdict. On CLINC150
the corresponding numbers are \EClincWThirtyTwo\% against \EClincWK\%, a
difference of \EClincDelta{} points and a lower bound of \EClincLower, so the
verdict is that the claim \EClincVerdict. On MASSIVE intent the width-32 arm
reaches \EMassiveWThirtyTwo\% against \EMassiveWK\%, a difference of
\EMassiveDelta{} points (lower bound \EMassiveLower), and the verdict is that the
claim \EMassiveVerdict. Across the three tasks the pre-registered claim
\EAggregateVerdict{} on \EHolds{} of \ETasks{} tasks.

\paragraph{The curve falls with width.} The non-inferiority test passes on every
task: the 32-way head is at or above the full bank everywhere, decisively so on
Banking77 and CLINC150 and within noise on MASSIVE intent. That is not a flat curve. Accuracy declines
monotonically as the candidate set widens. On Banking77, 2 candidates reach
\EBankWTwo\%, 8 reach \EBankWEight\%, 32 reach \EBankWThirtyTwo\% and the full bank
\EBankWK\%; CLINC150 gives \EClincWTwo, \EClincWEight, \EClincWThirtyTwo{} and
\EClincWK\%, and MASSIVE intent \EMassiveWTwo, \EMassiveWEight, \EMassiveWThirtyTwo{}
and \EMassiveWK\%. The width-8 arm is ahead of the width-32 arm by
\EBankWidthEightVsThirtyTwo{} points on Banking77, \EClincWidthEightVsThirtyTwo{} on CLINC150 and
\EMassiveWidthEightVsThirtyTwo{} on MASSIVE intent.

The arms are also not interchangeable. Re-scored on the full bank, the heads
trained at width 2 fall to \EBankFullKAtTwo, \EClincFullKAtTwo{} and
\EMassiveFullKAtTwo\%, and those trained at width 32 reach \EBankFullKAtThirtyTwo,
\EClincFullKAtThirtyTwo{} and \EMassiveFullKAtThirtyTwo\% on the full bank. Training at a
width teaches a head for that width; the at-width column is the comparison the
pre-registration fixed, and a narrow head is not a full-bank model that happens
to have seen fewer candidates.

So the registered claim survives, but the stronger reading does not. The 32-way
interface is at or above the full bank on every task, but because the full bank
is the hardest point in this curve, not because width is free. Widening from 2 to 32 costs \EBankCostTwoToThirtyTwo{} points on Banking77,
\EClincCostTwoToThirtyTwo{} on CLINC150 and \EMassiveCostTwoToThirtyTwo{} on MASSIVE intent. A product
that could pose 8 candidates would score higher than one posing 32 on these
three tasks; the 32-way interface is a compromise between coverage and accuracy,
not an accuracy-neutral choice. The valuation claim needs the registered
direction, and that direction is supported; it does not get the stronger claim
that accuracy is flat as the set widens.

The harness check is reported rather than assumed. The width-\(K\) arms score
\EBankWK, \EClincWK{} and \EMassiveWK\% against published suite values of
\EBankPublished, \EClincPublished{} and \EMassivePublished\%, so the absolute
deviations are \EBankHarnessDelta, \EClincHarnessDelta{} and
\EMassiveHarnessDelta{} points. The 0.2-point tolerance was set before the run
and two arms exceed it; the largest, \EMaxHarnessDelta{} points, is still inside
the \SuiteReplicationMax-point run-to-run floor the suite measured on repeated
cells. That is a level offset of the same size as the suite's own repeat noise,
and it is reported rather than hidden by substituting the published number.
The width comparison itself is within-run --- both arms trained by the same code
on the same rows --- so a level offset against the published cell shifts the
whole curve rather than the difference the rule tests.

What the sweep establishes is bounded. It is one seed per arm, one rung, three
English tasks and 4{,}000 training rows; the per-width arms differ in
optimisation trajectory as well as width, which is inherent to training the width
rather than simulating it. It says that at 113.47M a trained 32-way interface
is at or above the full bank on these cells, and that narrower interfaces do
better still;
it does not say that accuracy is constant in width, and it does not establish
near-perfect accuracy at any width. The remaining grid --- a second seed, and the
same twelve arms at the 14.4M and 33.7M rungs --- is costed in
\S\ref{sec:future}.
